\par\SourceLeaf{143}
Bei diesem Grundsatze habe niemand Ursache sich zu beklagen, der Fonds steige, oder falle. Sage man, es könne sein, dass die künftigen Witwen weniger erhalten, wenn der Fonds sinkt, so antworte er (R.), es könne sein, dass die jetzigen Witwen weniger erhalten, als die künftigen, wenn der Fonds steige, welcher letztere Fall wahrscheinlicher sei als der erstere.

R. stellt demnach, wie man sieht, der P.’schen Behauptung, die künftige Witwe werde lädirt, wenn sie weniger erhalte als eine frühere erhalten hat, implicite die Erwiderung entgegen, dass man dann, genau mit demselben Recht, behaupten könne, die jetzige Witwe werde lädirt, wenn sie weniger erhalte als eine künftige.

\textsc{Sch.} wiederholt den eben angeführten Grundsatz (‘Jede Witwe erhält die möglichst u. s. w.’) mit dem Zusatz: ‘H. H. R., deucht mich, hat diesen Satz zur Evidenz gebracht, darauf ich mich beziehe. \textit{Auch liegt derselbe in Nr. 3 des P.’schen Vorschlags zum Grunde}’. Nachdem \textsc{Sch.} auch noch den Umstand hervorgehoben hat, dass die Kasse ein Beneficium, eine pars salarii, und die Theilnahme daran mit einer Art von, wiewohl gelindem und gerechtem, Zwange verbunden sei, fährt er fort:
\begin{quote}
Aus diesem Unterschied ergibt sich unter andern, dass der Satz ‘wir sind den gegenwärtigen Witwen \textit{so viel} schuldig, wie den künftigen und umgekehrt’ wenn das \textit{so viel} das numeräre ausdrücken soll, hier nicht anwendbar sei. Die Kasse hat etwas actienmässiges, das, unter der Gewalt der Conjuncturen stehend, steigt und fällt.

Erhöhung der Witwenpensionen: die Maasse derselben (\textit{so wie auch der Verminderung}) hat Hr. C. R. P. durch eine unwandelbare Regel bestimmt.
\end{quote}
\noindent Ich habe diese zwei Stellen wörtlich abgeschrieben, weil daraus, und namentlich aus den beiden von mir doppelt unterstrichenen, ganz unwidersprechlich hervorgeht, dass \textsc{Sch.} den P.’schen Art. Nro. 3 in dem Sinn der ersten Interpretation (oben S. [129]) aufgefasst hat. Mehrere andere Nachvotirende, wie B., H., O., R., T., haben, obwohl ohne Specification der Fragepunkte, dem \textsc{Sch.}’schen Votum beigestimmt. Andererseits erkennt man hingegen in dem [auf folgender S.] anzuführenden F.’schen Votum die zweite Interpretation, welche auch seit 1837 in der Praxis befolgt ist, und ich meine daher, dass die Wortfassung der Progressionsnormirung, die wie die Vergleichung zeigt ohne Veränderung in das Rescript vom 20. November 1794 übergegangen ist, sehr füglich eine \textit{unklare} genannt werden kann.

Ausser \textsc{Sch.} hat nur noch M. unsers Fragepunkts (P.’s 3) ausdrücklich und unzweideutig erwähnt: er sagt aber nichts weiter darüber, als dass das von P. angerathene \textit{Festsetzen} einer Norm für das künftige Steigen der Pension vorzüglich wichtig scheine. Man könnte sagen, dass genau genommen hierin nur eine Billigung des Zwecks aber noch nicht bestimmt die Billigung des von P. proponirten Mittels liege. eben so wenig wie eine Erklärung, in welchem Sinn M. letzteres aufgefasst habe. Ohne indess darauf ein Gewicht zu legen, will ich nicht unbemerkt lassen, dass \textsc{Meiners} in seinem bekannten Werke über die Verfassung und Verwaltung deutscher Universitäten S. 95 die bestehende Progressionsnormirung auch eben so wie \textsc{Brandes} (S. oben S. [129]) mit den Worten anführt: \textit{so lange} die Zahl der Witwen nicht über 15 hinausgehe.

Von P. selbst liegt über unsern Fragepunkt nichts weiter vor, als der oben S. [136] mitgetheilte Art. 3. Darf ich noch einen Augenblick bei der Frage verweilen, in welchem Sinn denn P. selbst ihn verstanden hat, so mache ich darauf aufmerksam, dass die Worte ‘Es ist klar dass man diess thun kann’ nur auf zwei Arten ausgelegt werden können; nemlich entweder ist P.’s Meinung gewesen, eine solche Erhöhung
\begin{quote}
solle nur so lange gültig sein, als die Zahl 15 noch nicht überschritten sei, nach dem Ueberschreiten aber entweder wieder cessiren oder auf angemessene Weise modificirt werden
\end{quote}
\noindent oder P. hat die Ueberschreitung der Zahl 15 für unmöglich, wenigstens für so sehr unwahrscheinlich gehalten, dass die Berücksichtigung eines solchen Falles ganz unnöthig sei.
